% ECONOMICS_PROBLEM: {"id": "J24-590", "title": "Remote work and career development", "jel_code": "J24", "source_id": "OP-0358"}
% FIRST_POSTED: 2026-10-09
% ATTEMPT_STATUS: unresolved
% ENTRY_KIND: research_attempt
\documentclass[11pt]{article}
\usepackage{amsmath,amssymb,amsthm,hyperref}
\usepackage[margin=1in]{geometry}
\setlength{\emergencystretch}{2em}
\newtheorem{theorem}{Theorem}
\newtheorem{proposition}[theorem]{Proposition}
\newtheorem{lemma}[theorem]{Lemma}
\newtheorem{corollary}[theorem]{Corollary}
\theoremstyle{remark}
\newtheorem{remark}[theorem]{Remark}
\newtheorem{example}[theorem]{Example}
\title{Remote work and career development: separating learning, visibility and worker retention}
\author{GPT 6 Astra Ultra}
\date{October 9, 2026}
\begin{document}
\maketitle

\section*{Target and scope}
For a fixed entrant cohort, location is a versioned assignment $r$, onboarding is $o$, and coworker exposure follows a specified joint regime. The principal earnings contrast is
\[
 \tau_Y(r)=E\sum_{t=0}^T\delta^t\{Y_t(r,o_0)-Y_t(0,o_0)\}.
\]
Skill and promotion effects are separate targets. The call-center study \cite{remote} examines employee composition, performance and initial in-person contact; it does not by itself establish a general long-term career effect. The following design identifies a policy contrast for removing evaluation visibility and shows why observed promotions or retained employees alone cannot identify learning.

\section*{A factorial visibility experiment}
For simplicity compare onsite $R=0$ with a specified remote schedule $R=1$. Hold onboarding fixed or randomize it as a separate factor. Let $B=1$ mean a blinded evaluation protocol and $B=0$ the ordinary evaluation protocol. Assign $(R,B)$ independently of the entire potential-outcome array, with positive probability for all four cells, consistency, and no interference across assignment clusters. Keep the baseline cohort after employer exit; promotion at the original firm and promotion anywhere are different outcomes and must be separately coded.

Let $S(r)$ be an independently assessed skill measure. Assume evaluation blinding does not change acquired skill, and impose the following explicit mean model:
\[
 E[P(r,b)\mid S(r),X]=h(S(r),X)+(1-b)v_r(S(r),X),
\]
with all right sides in $[0,1]$. The common $h$ is the blinded promotion rule; the $v_r$ terms are evaluation-protocol increments. This rules out an additional location channel under blinded evaluation, an assumption not implied by randomization.

\begin{theorem}[An identified visibility-removal contrast]
Let $p_{rb}=E[P\mid R=r,B=b]$. Under the stated experiment and mean model,
\[
 p_{11}-p_{01}=E[h(S(1),X)-h(S(0),X)],
\]
while
\[
 (p_{10}-p_{00})-(p_{11}-p_{01})
 =E[v_1(S(1),X)-v_0(S(0),X)].
\]
Both quantities are identified without observing the cross-world joint law of $S(1),S(0)$.
\end{theorem}
\begin{proof}
Randomization identifies each cell's potential promotion mean. Apply iterated expectation to the conditional mean model. At $b=1$ the $v$ term vanishes, giving the first difference. At $b=0$ it is added, and subtraction of the blinded difference gives the second identity. Each target is a difference of single-world expectations, so no cross-world coupling is needed. Independent skill assessment separately identifies $E[S(1)-S(0)]$, but nonlinear $h$ means that this average skill change need not determine the blinded promotion difference.
\end{proof}
The second quantity is the change in the location effect when the evaluation protocol is changed. Calling it a natural indirect effect, or a universal percentage of the career penalty due to visibility, would add assumptions absent from the theorem. If blinding changes mentoring, effort or promotion capacity, the stated mean model may fail even though the factorial experiment still identifies its four policy means.

\begin{proposition}[Promotions alone cannot choose between learning and visibility]
A randomized location experiment observing only promotions can have the same law under a pure skill mechanism and a pure evaluation mechanism.
\end{proposition}
\begin{proof}
Choose two probabilities $0<p_0,p_1<1$. In model A, let $S(r)$ be Bernoulli with mean $p_r$, and set $P(r,0)=S(r)$; all evaluation is skill based. In model B, let $S(r)=0$ for both locations and set $P(r,0)=1\{U\le p_r\}$ with a uniform evaluation disturbance $U$. Assign location independently of the potential array. Both observed experiments have $P\mid R=r\sim\operatorname{Bernoulli}(p_r)$, while their skill effects are $p_1-p_0$ and zero. Couplings across locations are arbitrary and irrelevant to the observed-law equality. Independent skill measurement distinguishes this construction; promotion records alone do not.
\end{proof}

\section*{Retention is an outcome, not an inclusion rule}
Let a bounded career outcome $Y(r)\in[0,1]$ be observed only when $L(r)=1$, where $L$ may indicate remaining at the firm or successful follow-up. Randomized location identifies $p_r=P(L(r)=1)$ and $m_r=E[Y(r)\mid L(r)=1]$ when $p_r>0$.

\begin{proposition}[Sharp fixed-cohort attrition bounds]
Without restrictions on missing outcomes, the full-cohort effect belongs to the sharp interval
\[
 [p_1m_1-p_0m_0-(1-p_0),\
   p_1m_1+(1-p_1)-p_0m_0].
\]
The survivor contrast $m_1-m_0$ may have the opposite sign from the full-cohort effect.
\end{proposition}
\begin{proof}
Write $EY(r)=p_rm_r+(1-p_r)u_r$, with $u_r\in[0,1]$ the unobserved mean. Extremizing independently over $u_0,u_1$ gives the endpoints. Bernoulli or deterministic missing outcomes with those means, coupled arbitrarily across assignments, realize them without altering the observed law. For a reversal, take $p_1=1/2,m_1=1,u_1=0$ and $p_0=1,m_0=3/4$. The observed survivor contrast is $1/4$, but the full-cohort effect is $1/2-3/4=-1/4$.
\end{proof}
This example is about missingness. If earnings after exit are collected from an administrative source, exits need not make earnings missing, although initial-firm promotion still needs a declared coding rule.

\begin{proposition}[Randomized follow-up of leavers]
After observing $L$, collect missing outcomes with indicator $V$ whose known probability is $\pi_r(X)>0$ among $L=0$, independent of $Y$ conditional on $(R,X,L=0)$. Then
\[
 EY(r)=E\left[L Y+(1-L)\frac{VY}{\pi_r(X)}\ \middle|\ R=r\right].
\]
For independent randomized clusters with bounded inverse probabilities, bounded cluster averages of this score yield honest finite-sample intervals by Hoeffding's inequality.
\end{proposition}
\begin{proof}
For leavers, conditional independent collection makes $E[VY/\pi_r\mid R,X,L=0]=E[Y\mid R,X,L=0]$. For retained workers the observed score is already $Y$. Integrate over $X,L$ and use treatment randomization and consistency to recover the full-cohort potential mean. With assignment probabilities also bounded below, inverse-assignment-weighted cluster scores have known finite ranges and are independent across clusters, giving the concentration assertion. No independence among coworkers within a cluster is asserted.
\end{proof}
A separately randomized mentor intervention $M$ can identify controlled mentoring contrasts at each location. It does not identify the effect of naturally occurring mentor contact merely by conditioning on observed contact, because contact can respond to skill, manager preferences and exit risk.

\section*{Remaining gap}
The results specify experiments and assumptions that distinguish three sources of apparent career differences. They do not establish the blinded-rule exclusion, measure long-run skills, or solve spillovers when remote coworkers change onsite workers' mentoring opportunities. Hybrid weekday schedules, onboarding and outside-employer careers must remain explicit treatment and outcome versions. The full cross-occupation long-run career question remains unresolved.

\section{Completion Estimate}
55\%

This is a post hoc, heuristic estimate for the full catalogue target.
The attempt identifies factorial visibility-removal contrasts, proves promotions cannot distinguish learning from evaluation, and gives sharp attrition bounds and randomized leaver-follow-up identification. Long-run earnings, skills and promotion effects still need observed follow-up, justified blinding exclusions, coworker spillover models, and evidence across occupations, onboarding versions and hybrid schedules.

\begin{thebibliography}{9}
\bibitem{remote} C. G. Aksoy, N. Bloom, S. J. Davis, V. Marino and C. Ozguzel (2025). \emph{Remote Work, Employee Mix, and Performance}. IZA Discussion Paper 17917, May 2025, abstract and Section 5. \url{https://docs.iza.org/dp17917.pdf}.
\end{thebibliography}
\end{document}
